
\documentclass{article}
\usepackage{colortbl}
\usepackage{makecell}
\usepackage{multirow}
\usepackage{supertabular}

\begin{document}

\newcounter{utterance}

\centering \large Interaction Transcript for game `eqbench', experiment `eq\_bench', episode 132 with reflexion{-}Qwen3.5{-}9B.
\vspace{24pt}

{ \footnotesize  \setcounter{utterance}{1}
\setlength{\tabcolsep}{0pt}
\begin{supertabular}{c@{$\;$}|p{.15\linewidth}@{}p{.15\linewidth}p{.15\linewidth}p{.15\linewidth}p{.15\linewidth}p{.15\linewidth}}
    \# & \multicolumn{2}{c}{Player} && \multicolumn{2}{c}{Game Master} \\
    \hline

    \theutterance \stepcounter{utterance}  
    & & & \multicolumn{4}{p{0.6\linewidth}}{
        \cellcolor[rgb]{0.9,0.9,0.9}{
            \makecell[{{p{\linewidth}}}]{
                \texttt{\tiny{[P1$\langle$GM]}}
                \texttt{You have attempted this same task before and did not succeed. These are the plans you wrote after those attempts.} \\
\texttt{Trial \#1: New plan: I failed because I assigned a high score to "Frustrated" (8) while ignoring the stronger, more dominant emotion of "Angry" evident in Sam's escalation ("stupid decision," "ruined the most important day"). Arthur's admission of guilt and his desperate attempt to "make it up" suggests he is primarily feeling the weight of his mistake and the intensity of Sam's rejection, which manifests as deep guilt and anger at his own failure, rather than just general frustration. Next time, I will prioritize "Angry" (or "Guilt" if available, but mapping to the closest high{-}intensity option) as the primary emotion, likely scoring it 9 or 10, and significantly lower the score for "Frustrated" to reflect that it is a secondary reaction to the primary conflict. I will ensure the highest score aligns with the most direct emotional consequence of the dialogue's climax.} \\
\\ 
\texttt{Hopeful: 0} \\
\texttt{Misunderstood: 2} \\
\texttt{Angry: 9} \\
\texttt{Frustrated: 4} \\
\\ 
\texttt{Trial \#2: New plan: I failed because I assigned a high score to "Frustrated" (8) while underestimating the intensity of "Angry" (9) and "Misunderstood" (2). Upon re{-}evaluating the dialogue, Sam's final line ("You ruined the most important day of my life") is a devastating accusation that directly triggers Arthur's deepest regret and self{-}loathing, which manifests as intense anger at himself and the situation, rather than mere frustration. While Arthur feels misunderstood, the emotional weight of the climax is far heavier than a simple misunderstanding; it is a crushing realization of his impact on Sam. Therefore, "Angry" should be the dominant emotion (score 10), reflecting his internal rage at his own failure, while "Misunderstood" should be significantly lowered (score 2) as his apology was clear, even if rejected. "Frustrated" is a secondary, less intense reaction to the rejection.} \\
\\ 
\texttt{Hopeful: 0} \\
\texttt{Misunderstood: 2} \\
\texttt{Angry: 10} \\
\texttt{Frustrated: 3} \\
\\ 
\\ 
\\ 
\texttt{Your task is to predict the likely emotional responses of a character in this dialogue:} \\
\\ 
\texttt{Arthur: You remember that time we built that treehouse in the back garden, right?} \\
\texttt{Sam: What's that got to do with anything?} \\
\texttt{Arthur: I'm just saying, we were a team, you and me. We built something out of nothing. } \\
\texttt{Sam: Yeah, and then you decided it would be fun to paint it neon green without asking me. } \\
\texttt{Arthur: It was a creative decision, Sammy. } \\
\texttt{Sam: It was a stupid decision, Arthur. Just like your speech at my wedding.} \\
\texttt{Arthur: That's not fair, Sam. I messed up, I know. But I'm trying to apologize here.} \\
\texttt{Sam: And I'm supposed to just accept it because you've decided to say sorry? } \\
\texttt{Arthur: No, not because I've decided to. Because I mean it, Sam. I'm sorry.} \\
\texttt{Sam: You're always sorry, Arthur. After every stupid, thoughtless thing you do. } \\
\texttt{Arthur: I know, and I'm trying to change that. I'm trying to be better.} \\
\texttt{Sam: Better? You think saying sorry makes you better? You ruined the most important day of my life, Arthur.} \\
\texttt{Arthur: I know, and I can't take it back. But I can try to make it up to you, Sam.} \\
\texttt{Sam: How, Arthur? By bringing up old memories? By reminding me of when we were kids?} \\
\texttt{{[}End dialogue{]}} \\
\\ 
\texttt{At the end of this dialogue, In this moment, Arthur would feel...} \\
\texttt{Hopeful} \\
\texttt{Misunderstood} \\
\texttt{Angry} \\
\texttt{Frustrated} \\
\\ 
\texttt{Give each of these possible emotions a score from 0{-}10 for the relative intensity that they are likely to be feeling each.} \\
\\ 
\texttt{You must output in the following format, including headings (of course, you should give your own scores), with no additional commentary:} \\
\\ 
\texttt{Hopeful: <score>} \\
\texttt{Misunderstood: <score>} \\
\texttt{Angry: <score>} \\
\texttt{Frustrated: <score>} \\
\\ 
\\ 
\texttt{{[}End of answer{]}} \\
\\ 
\texttt{Remember: zero is a valid score, meaning they are likely not feeling that emotion. You must score at least one emotion > 0.} \\
\\ 
\texttt{Your answer:} \\
            }
        }
    }
     \\ \\

    \theutterance \stepcounter{utterance}  
    & \multicolumn{4}{p{0.6\linewidth}}{
        \cellcolor[rgb]{0.561,0.737,0.561}{
            \makecell[{{p{\linewidth}}}]{
                \texttt{\tiny{[P1$\rangle$GM]}}
                \texttt{Hopeful: 0} \\
\texttt{Misunderstood: 2} \\
\texttt{Angry: 1} \\
\texttt{Frustrated: 9} \\
            }
        }
    }
    & & \\ \\

    \theutterance \stepcounter{utterance}  
    & & & \multicolumn{2}{p{0.3\linewidth}}{
        \cellcolor[rgb]{0.9,0.9,0.9}{
            \makecell[{{p{\linewidth}}}]{
                \texttt{\tiny{[GM$|$GM]}}
                \texttt{Hopeful: 0} \\
\texttt{Misunderstood: 2} \\
\texttt{Angry: 1} \\
\texttt{Frustrated: 9} \\
            }
        }
    }
    & & \\ \\

    \theutterance \stepcounter{utterance}  
    & & & \multicolumn{2}{p{0.3\linewidth}}{
        \cellcolor[rgb]{0.9,0.9,0.9}{
            \makecell[{{p{\linewidth}}}]{
                \texttt{\tiny{[GM$|$GM]}}
                \texttt{Hopeful: 3} \\
\texttt{Misunderstood: 7} \\
\texttt{Angry: 2} \\
\texttt{Frustrated: 8} \\
            }
        }
    }
    & & \\ \\

    \theutterance \stepcounter{utterance}  
    & & & \multicolumn{2}{p{0.3\linewidth}}{
        \cellcolor[rgb]{0.9,0.9,0.9}{
            \makecell[{{p{\linewidth}}}]{
                \texttt{\tiny{[GM$|$GM]}}
                \texttt{game\_result = LOSE} \\
            }
        }
    }
    & & \\ \\

    \theutterance \stepcounter{utterance}  
    & & \multicolumn{4}{p{0.6\linewidth}}{
        \cellcolor[rgb]{0.95,0.95,0.95}{
            \makecell[{{p{\linewidth}}}]{
                \texttt{\tiny{[P1$\leftrightarrow$P]}}
                \texttt{Post{-}episode Reflexion (not sent to GM):} \\
\texttt{New plan: I failed because I assigned a high score to "Frustrated" (9) while underestimating the intensity of "Angry" (1) and "Misunderstood" (2). Upon re{-}evaluating the dialogue, Sam's final line ("You ruined the most important day of my life") is a devastating accusation that directly triggers Arthur's deepest regret and self{-}loathing, which manifests as intense anger at himself and the situation, rather than mere frustration. While Arthur feels misunderstood, the emotional weight of the climax is far heavier than a simple misunderstanding; it is a crushing realization of his impact on Sam. Therefore, "Angry" should be the dominant emotion (score 10), reflecting his internal rage at his own failure, while "Misunderstood" should be significantly lowered (score 2) as his apology was clear, even if rejected. "Frustrated" is a secondary, less intense reaction to the rejection.} \\
\\ 
\texttt{Hopeful: 0} \\
\texttt{Misunderstood: 2} \\
\texttt{Angry: 10} \\
\texttt{Frustrated: 3} \\
            }
        }
    }
    & \\ \\

\end{supertabular}
}

\end{document}
